% ==============================================================================
%  1.1  Modelos estadísticos: utilidad y fundamentos
% ==============================================================================
\section{Modelos estadísticos: utilidad y fundamentos}
\label{sec:t1-modelos}

\subsection{La tarea del estadístico}

La tarea de un estadístico es \textbf{extraer información de un conjunto de datos} para:
\begin{itemize}
  \item describir un fenómeno,
  \item controlar o predecir el fenómeno,
  \item confirmar alguna hipótesis sobre el fenómeno,
  \item informar para la toma de decisiones.
\end{itemize}

\paragraph{El proceso estadístico.} De una \emph{población} se extrae una \emph{muestra}
(muestreo), que se resume mediante la \emph{estadística descriptiva}. Para extraer
conclusiones sobre la población a partir de la muestra (\emph{inferencia}) se necesita la
\emph{probabilidad}, es decir, describir cómo se han generado los datos de la muestra. Esta es
la tarea de la \textbf{modelización estadística}.

\begin{figure}[H]
  \centering
  \begin{tikzpicture}[
      caja/.style={draw, rounded corners, fill=fondoDef, minimum width=2.6cm,
                   minimum height=0.9cm, align=center, font=\small},
      flecha/.style={-{Stealth[length=2.5mm]}, thick}]
    \node[caja] (pob) at (0,0) {Población};
    \node[caja] (mue) at (7,0) {Muestra};
    \draw[flecha] (pob.north east) to[bend left=25]
      node[above, font=\small] {Muestreo} (mue.north west);
    \node[font=\small, right=0.3cm of mue] {$\to$ Estadística descriptiva};
    \draw[flecha, colRes!80!black] (mue.south west) to[bend left=25]
      node[below, font=\small, align=center]
      {\textbf{Probabilidad} + \textbf{Inferencia}\\[-1pt]
       \footnotesize (modelización estadística)} (pob.south east);
  \end{tikzpicture}
  \caption{El proceso estadístico. La modelización estadística permite pasar de la muestra
    a la población.}
\end{figure}

\subsection{Componente sistemática y componente aleatoria}

Los datos suelen combinar dos aspectos:
\begin{itemize}
  \item cierta regularidad en algún sentido: \textbf{componente sistemática};
  \item una variación irregular: \textbf{componente aleatoria}.
\end{itemize}

\begin{ejemplo}[Descomposición sistemática + aleatoria][ej:t1-componentes]
  \begin{enumerate}[label=(\alph*)]
    \item \emph{Consumo de cigarrillos de un fumador.} Varía irregularmente de semana en
      semana en torno a un nivel ``medio'': $C_i = m + e_i$.
    \item \emph{Ventas de patatas en un establecimiento.} Oscilan a lo largo del año con un
      patrón estacional, pero no todas las semanas son iguales: $V_t = m(t) + e_t$.
    \item \emph{Elección entre dos rutas, A y B, para cruzar la ciudad.} Los datos son los
      tiempos de viaje durante varios días usando las dos rutas. El tiempo de cada viaje se
      descompone en una componente fija, la ``norma'', y otra aleatoria (tráfico,
      semáforos\ldots): $t_r = T + u_r$, $r=1,\ldots,n$.
  \end{enumerate}
  En los tres casos, $m$, $m(t)$ y $T$ son la componente sistemática y $e_i$, $e_t$, $u_r$ la
  aleatoria.

  En (c), para saber qué ruta es más rápida no basta con comparar tiempos individuales, ya que
  cada uno incluye la componente aleatoria $u_r$. Hay que comparar las componentes fijas $T_A$
  y $T_B$ de cada ruta, lo que supone hacer inferencia sobre la diferencia de dos medias a
  partir de las dos muestras de tiempos (\cref{sec:t1-dos-medias}).
\end{ejemplo}

\subsection{Modelo estadístico}

\begin{definicion}[Modelo estadístico]
  La característica principal de un modelo estadístico es que \textbf{la variabilidad se
  representa por medio de distribuciones de probabilidad}. Por ello, la formulación del
  modelo incluye alguna especificación sobre la componente aleatoria (o \emph{perturbación}),
  como que sigue cierta distribución de probabilidad, por ejemplo $\Normal(0,\sigma^2)$.
\end{definicion}

\paragraph{Etapas del ajuste de un modelo estadístico.}
\begin{enumerate}
  \item \textbf{Especificación}: seleccionar una función apropiada para el modelo y las
    variables a incluir. Puede ser muy útil el análisis exploratorio de los datos
    (\cref{sec:t1-descriptivo}). Para un mismo conjunto de datos pueden ser apropiados
    diferentes modelos y tipos de análisis.
  \item \textbf{Estimación--ajuste}: seleccionar un método apropiado de estimación del
    modelo y estimar los parámetros.
  \item \textbf{Diagnóstico}: comprobar si el modelo propuesto es correcto. A veces esta etapa sirve
    para seleccionar ``el mejor modelo''.
\end{enumerate}

\paragraph{Tipos de modelos estadísticos.}
\begin{itemize}
  \item \textbf{Paramétricos}: el modelo queda perfectamente especificado una vez conocido
    el conjunto de parámetros que, a su vez, especifican la distribución de probabilidad
    subyacente. \emph{Ejemplo}: si la componente aleatoria tiene distribución
    $\Normal(0,\sigma^2)$, para estimar el modelo hay que estimar $\sigma$.
  \item \textbf{No paramétricos}: la distribución de los datos no puede definirse en
    términos de un conjunto de parámetros. \emph{Ejemplo}: se suponen las perturbaciones
    independientes y con media 0, sin asumir ninguna distribución concreta.
\end{itemize}

\begin{nota}[Notación]
  En estos apuntes $\Normal(\mu,\sigma^2)$ indica siempre una normal de media $\mu$ y
  \emph{varianza} $\sigma^2$ (desviación típica $\sigma$).
\end{nota}

\subsection{El modelo de regresión}

En muchos problemas existe una relación entre dos o más variables cuya naturaleza interesa
estudiar. El \textbf{análisis de regresión} es la técnica estadística que se utiliza para
modelizar e investigar este tipo de relaciones.

\begin{definicion}[Modelo y función de regresión][def:funcion-regresion]
  Los \textbf{modelos de regresión} describen cómo una \emph{variable respuesta} $Y$, tratada
  como \emph{aleatoria}, depende de un conjunto de \emph{variables explicativas}
  $X_1,\ldots,X_k$, tratadas como \emph{fijas}. Se supone que $Y$ tiene una distribución de
  probabilidad para cada nivel de las $X$, cuya media varía con ellas de forma sistemática.
  Esta relación sistemática es la \textbf{función de regresión de $Y$ sobre $X$}:
  \[
    m(x_1,\ldots,x_k) = \E\left(Y \cond X_1=x_1,\ldots,X_k=x_k\right).
  \]
\end{definicion}

Cada forma de la función de regresión y de la distribución de $Y$ da lugar a un modelo de
regresión distinto, que se estima a partir de los datos. Cuando no se hace ninguna suposición
sobre la forma de $m$ se obtienen los \textbf{modelos de regresión no paramétricos}, que se
ajustan estimando $m(x_1,\ldots,x_k)$ punto a punto.

\subsection{El modelo lineal}

\begin{definicion}[Modelo lineal][def:modelo-lineal]
  Un modelo es \textbf{lineal} si su componente sistemática es una \textbf{función lineal de
  sus parámetros}. Los ejemplos más sencillos son la constante $\alpha$ y la recta
  $\alpha+\beta x$. En ningún caso aparecen expresiones del tipo $\alpha^2$, $\alpha\beta$,
  $\log(\beta)$\ldots
\end{definicion}

\begin{nota}
  ``Lineal'' se refiere a los \emph{parámetros}, no a las variables: $Y=\beta_0+\beta_1X^2+U$
  o $Y=\beta_0+\beta_1\log X+U$ son modelos lineales; $Y=\beta_0 e^{\beta_1X}+U$ no lo es.
\end{nota}

Formalmente, el \textbf{modelo lineal general} se expresa como
\[
  \vect{Y} = Z\vect{\theta} + \vect{U},
\]
donde $\vect{Y}$ es el vector de la variable respuesta (univariante o multivariante), $Z$ una
matriz conocida que contiene los valores de las variables explicativas, $\vect{\theta}$ un
vector de parámetros y $\vect{U}$ un vector aleatorio. Puede acompañarse de un conjunto de
restricciones lineales sobre los parámetros, $R\vect{\theta}=\vect{c}$.

\paragraph{Modelos de regresión lineal.} Son aquellos en los que una respuesta continua depende
linealmente de un conjunto de variables explicativas,
\[
  Y = \beta_0 + \beta_1X_1 + \cdots + \beta_kX_k + U .
\]
Constituyen la base de la mayoría de los análisis estadísticos, aunque no son adecuados en
todos los casos: en muchas aplicaciones la respuesta $Y$ es \emph{discreta}, o los datos
sugieren la presencia de \emph{componentes no lineales}. En esas situaciones se recurre a:
\begin{itemize}
  \item \textbf{modelos no lineales} $Y=g(X_1,\ldots,X_k)+U$ con $g$ no lineal, que requieren
    nuevas herramientas de ajuste (por ejemplo, mínimos cuadrados ponderados);
  \item \textbf{modelos lineales generalizados}, con distribuciones de la familia
    exponencial para la respuesta, de la que la binomial y la Poisson son casos particulares
    importantes.
\end{itemize}
